\documentclass[12pt,a4paper]{article}

% Pacotes principais
\usepackage{amssymb, amsmath, graphicx, float, enumitem, caption, hyperref, fancyhdr}
\usepackage[margin=2.5cm]{geometry} % margens ajustadas
\usepackage[utf8]{inputenc}
\usepackage[T1]{fontenc}
\usepackage{lmodern} % fonte mais limpa e moderna
\usepackage{setspace} % controle de espaçamento
\usepackage{titlesec} % customizar títulos
\usepackage{xcolor}   % cores opcionais

% Espaçamento entre linhas
\onehalfspacing

% Formato dos títulos (corrigido)
\titleformat{\section}{\bfseries\Large}{\thesection}{1em}{}
\titleformat{\subsection}{\bfseries\large}{\thesubsection}{0.5em}{}

% Espaçamento entre listas
\setlist[itemize]{itemsep=0.5em, topsep=0.5em}

% Cabeçalho e rodapé
\pagestyle{fancy}
\fancyhf{}
\fancyhead[L]{Lista de Exercícios}
\fancyhead[R]{João Pedro Rolim Ximenes}
\fancyfoot[C]{\thepage}

\title{\vspace{-2cm}Lista de Exercícios\vspace{-0.5cm}}
\author{João Pedro Rolim Ximenes \\ 566235}
\date{\vspace{-1.5cm}}

\begin{document}

% Capa
\begin{center}
    \includegraphics[scale=0.9]{imagens/ufc.png} \\[1cm]
    \textbf{\LARGE Universidade Federal do Ceará} \\[2cm]
    \textbf{\huge EXERCÍCIOS - AULAS 1 A 3} \\[3cm]
    \textbf{\Large João Pedro Rolim Ximenes} \\[0.2cm]
    566235
\end{center}

\newpage

\section*{AULA 01}

\subsection*{Exercício 1}
Classifique as variáveis abaixo:
\begin{itemize}
    \item Preço do kg de pão em Reais - Quantitativo
    \item Renda familiar - Quantitativo
    \item Salário em Reais - Quantitativo
    \item Classe de renda  - Qualitativo
    \item Horas de estudo por dia - Quantitativo
    \item Tempo de estudo diário - Qualitativo
    \item Nível de emprego - Qualitativo
    \item Atividade profissional - Qualitativo
    \item Nível de satisfação do cliente - Qualitativo
    \item \textit{Explicação:} Se for baseado em contagens objetivas é quantitativo; se for baseado em avaliações de qualidade é qualitativo.
\end{itemize}

\subsection*{Exercício 2}
Dados: $X = 2, 1, 2, 1, 3, 2, 0, 2, 0, 1, 2, 2$

\begin{itemize}
    \item Classificação da variável: Quantitativa discreta
    \item Número de unidades e níveis: 12 unidades e 4 níveis
\end{itemize}

\begin{table}[H]
\centering
\caption*{Frequência Absoluta}
\renewcommand{\arraystretch}{1.2}
\begin{tabular}{|c|c|}
\hline
Valor & Frequência \\ \hline
0 & 2 \\ \hline
1 & 3 \\ \hline
2 & 6 \\ \hline
3 & 1 \\ \hline
\end{tabular}
\end{table}

\begin{table}[H]
\centering
\caption*{Frequência Relativa}
\renewcommand{\arraystretch}{1.2}
\begin{tabular}{|c|c|}
\hline
Valor & Frequência Relativa \\ \hline
0 & 0,17 \\ \hline
1 & 0,25 \\ \hline
2 & 0,50 \\ \hline
3 & 0,08 \\ \hline
\end{tabular}
\end{table}

\begin{table}[H]
\centering
\caption*{Frequência Percentual}
\renewcommand{\arraystretch}{1.2}
\begin{tabular}{|c|c|}
\hline
Valor & Percentual \\ \hline
0 & 17\% \\ \hline
1 & 25\% \\ \hline
2 & 50\% \\ \hline
3 & 8\% \\ \hline
\end{tabular}
\end{table}

\subsection*{Exercício 3}
\begin{itemize}
    \item O número de unidades estatísticas menores que 20 anos: 
    \[
    (0,1 \times 120) + (0,2 \times 80) = 28 \text{ unidades}
    \]
    \item O percentual de unidades com idade maior ou igual a 50 anos: 
    \[
    \big((0,5 \times 120) + (0,3 \times 80)\big) \times 0,5 = 42\%
    \]
    \item O número de homens com 30 anos ou mais: 
    \[
    0,8 \times 120 = 96 \text{ homens}
    \]
\end{itemize}

\section*{AULA 03}

\subsection*{Exercício 2}
\begin{figure}[H]
    \centering
    \includegraphics[width=0.7\linewidth]{imagens/boxplot.jpg}
    \caption*{Box-Plot}
\end{figure}

\subsection*{Exercício 3}
\begin{figure}[H]
    \centering
    \includegraphics[width=0.7\linewidth]{imagens/boxplot2.jpg}
\end{figure}

\begin{figure}[H]
    \centering
    \includegraphics[width=0.7\linewidth]{imagens/grafico linha.jpg}
\end{figure}

\begin{itemize}
    \item Mediana: 20
    \item Variância: 4,25
    \item Desvio padrão: 2,06
    \item Mínimo: 14
    \item Q1: 19
    \item Q2: 20
    \item Q3: 21
    \item Máximo: 25
\end{itemize}

\end{document}
